% Slide 3: Regulatory Alignment & Technical Context
\begin{frame}
  \frametitle{From Academic Critique to Legal Compliance}
  \framesubtitle{Regulatory Imperative \textbar\ Lead Presenter: Tisha}

  \begin{columns}[T]
    \begin{column}{0.48\textwidth}
      \begin{card}{EU AI Act (Regulation EU 2024/1689)}
        \scriptsize
        \textbf{Phased Applicability via Article 113:}
        \begin{itemize}\scriptsize
          \item \textbf{Article 10 (Data Governance):} Datasets must be examined for \textit{"possible biases that are likely to lead to discrimination."}
          \item \textbf{Annex III Classification:} Biometric categorization inferring sensitive traits (race, age, gender) is classified as \hlcoral{High-Risk AI}.
          \item \textbf{Article 15 (Accuracy \& Robustness):} Requires declared demographic performance metrics.
        \end{itemize}
      \end{card}
    \end{column}

    \begin{column}{0.48\textwidth}
      \begin{card}{NIST AI RMF 1.0 (Measure Function)}
        \scriptsize
        \textbf{Voluntary Measurement Taxonomy:}
        \begin{itemize}\scriptsize
          \item \textbf{Measure 2.11:} Systematic measurement of bias and demographic parity.
          \item \textbf{Measure 1.1:} Documenting unquantifiable or low-sample risks ($n < 30$ guard).
          \item \textbf{Measure 1.3:} Software implementation cross-checking.
        \end{itemize}
      \end{card}

      \vspace{0.1cm}
      \begin{statcard}{The Technical Audit Mandate}
        \scriptsize Compliance requires \textbf{reproducible, offline-auditable software infrastructure} translating legal mandates into concrete statistical metrics.
      \end{statcard}
    \end{column}
  \end{columns}
\end{frame}
